
\documentclass{article}
\usepackage{colortbl}
\usepackage{makecell}
\usepackage{multirow}
\usepackage{supertabular}

\begin{document}

\newcounter{utterance}

\centering \large Interaction Transcript for game `wordle', experiment `high\_frequency\_words\_no\_clue\_no\_critic', episode 4 with deepseek{-}v4{-}pro.
\vspace{24pt}

{ \footnotesize  \setcounter{utterance}{1}
\setlength{\tabcolsep}{0pt}
\begin{supertabular}{c@{$\;$}|p{.15\linewidth}@{}p{.15\linewidth}p{.15\linewidth}p{.15\linewidth}p{.15\linewidth}p{.15\linewidth}}
    \# & \multicolumn{2}{c}{Player} && \multicolumn{2}{c}{Game Master} \\
    \hline

    \theutterance \stepcounter{utterance}  
    & & & \multicolumn{4}{p{0.6\linewidth}}{
        \cellcolor[rgb]{0.9,0.9,0.9}{
            \makecell[{{p{\linewidth}}}]{
                \texttt{\tiny{[P1$\langle$GM]}}
                \texttt{Jste jazykový kouzelník, který rád hádá slova podle daných pravidel.} \\
\\ 
\texttt{Vítejte ve hře Wordle! Máte šest pokusů na uhodnutí cílového slova, platného českého slova o pěti malých písmenech. Použijte značky "vysvětlení:" a "hádání:" k poskytnutí stručného vysvětlení pro každý odhad.} \\
\\ 
\texttt{Například, pokud je váš odhad "jablko", vaše odpověď by měla být} \\
\texttt{vysvětlení: toto je běžné české slovo o pěti písmenech a začínám svůj odhad tímto slovem.} \\
\texttt{hádání: jablko} \\
\\ 
\texttt{Po každém odhadu bude vaše odpověď ověřena a obdržíte zpětnou vazbu označující, která písmena jsou správná (zelená), která písmena jsou správná, ale na špatné pozici (žlutá), a která písmena jsou nesprávná (červená). Tato zpětná vazba může být užitečná při určování, která písmena zahrnout nebo vyloučit v dalším odhadu.} \\
\\ 
\texttt{Například zpětná vazba pro "jablko" může být:} \\
\texttt{zpětná vazba k hádání: d<žlutá> o<žlutá> l<zelená> a<žlutá> r<červená> o<červená>} \\
\\ 
\texttt{Vysvětlení by mělo obsahovat podrobnosti o tom, jak je zpětná vazba k hádání použita k vytvoření nového odhadu.} \\
\\ 
\texttt{Nevytvářejte žádný jiný text.} \\
\\ 
\texttt{Začněme vaším prvním odhadem.} \\
            }
        }
    }
     \\ \\

    \theutterance \stepcounter{utterance}  
    & \multicolumn{4}{p{0.6\linewidth}}{
        \cellcolor[rgb]{0.561,0.737,0.561}{
            \makecell[{{p{\linewidth}}}]{
                \texttt{\tiny{[P1$\rangle$GM]}}
                \texttt{vysvětlení: začínám běžným českým slovem o pěti písmenech, které obsahuje více samohlásek a časté souhlásky pro získání základní informace.} \\
\texttt{hádání: strom} \\
            }
        }
    }
    & & \\ \\

    \theutterance \stepcounter{utterance}  
    & & & \multicolumn{2}{p{0.3\linewidth}}{
        \cellcolor[rgb]{0.9,0.9,0.9}{
            \makecell[{{p{\linewidth}}}]{
                \texttt{\tiny{[GM$|$GM]}}
                \texttt{attempts = 1} \\
\texttt{target\_word = sídlo} \\
\texttt{guess = strom} \\
\texttt{guess\_feedback = s<green> t<red> r<red> o<yellow> m<red>} \\
            }
        }
    }
    & & \\ \\

    \theutterance \stepcounter{utterance}  
    & & & \multicolumn{4}{p{0.6\linewidth}}{
        \cellcolor[rgb]{0.9,0.9,0.9}{
            \makecell[{{p{\linewidth}}}]{
                \texttt{\tiny{[P1$\langle$GM]}}
                \texttt{zpětná vazba k hádání: s<zelená> t<červená> r<červená> o<žlutá> m<červená>} \\
\\ 
\texttt{Odpovězte pouze v tomto formátu.} \\
\texttt{vysvětlení: podrobnosti} \\
\texttt{hádání: slovo} \\
            }
        }
    }
     \\ \\

    \theutterance \stepcounter{utterance}  
    & \multicolumn{4}{p{0.6\linewidth}}{
        \cellcolor[rgb]{0.561,0.737,0.561}{
            \makecell[{{p{\linewidth}}}]{
                \texttt{\tiny{[P1$\rangle$GM]}}
                \texttt{vysvětlení: písmeno 's' je zelené (správné na první pozici), 'o' je žluté (správné, ale na jiné pozici). 't', 'r', 'm' jsou červené (nejsou ve slově). hledám slovo začínající na 's' obsahující 'o' na jiné pozici, bez 't', 'r', 'm'. zkouším 'sokol'.} \\
\texttt{hádání: sokol} \\
            }
        }
    }
    & & \\ \\

    \theutterance \stepcounter{utterance}  
    & & & \multicolumn{2}{p{0.3\linewidth}}{
        \cellcolor[rgb]{0.9,0.9,0.9}{
            \makecell[{{p{\linewidth}}}]{
                \texttt{\tiny{[GM$|$GM]}}
                \texttt{attempts = 2} \\
\texttt{target\_word = sídlo} \\
\texttt{guess = sokol} \\
\texttt{guess\_feedback = s<green> o<yellow> k<red> o<red> l<yellow>} \\
            }
        }
    }
    & & \\ \\

    \theutterance \stepcounter{utterance}  
    & & & \multicolumn{4}{p{0.6\linewidth}}{
        \cellcolor[rgb]{0.9,0.9,0.9}{
            \makecell[{{p{\linewidth}}}]{
                \texttt{\tiny{[P1$\langle$GM]}}
                \texttt{zpětná vazba k hádání: s<zelená> o<žlutá> k<červená> o<červená> l<žlutá>} \\
\\ 
\texttt{Odpovězte pouze v tomto formátu.} \\
\texttt{vysvětlení: podrobnosti} \\
\texttt{hádání: slovo} \\
            }
        }
    }
     \\ \\

    \theutterance \stepcounter{utterance}  
    & \multicolumn{4}{p{0.6\linewidth}}{
        \cellcolor[rgb]{0.561,0.737,0.561}{
            \makecell[{{p{\linewidth}}}]{
                \texttt{\tiny{[P1$\rangle$GM]}}
                \texttt{vysvětlení: první pozice je 's' (zelená). 'o' je žluté (bylo ve slově na druhé pozici, ale teď je červené, takže 'o' nemůže být na druhé ani čtvrté pozici; zůstává žluté z prvního pokusu). 'l' je žluté (nové písmeno). 'k' je červené. vyloučeno: t, r, m, k. známá písmena: s (1.), o (někde jinde), l (někde jinde). hledám slovo: s \_ \_ \_ \_. možnosti: slovo obsahuje 'o' a 'l', ne na pozicích 2 a 4 pro 'o', ne na pozici 5 pro 'l'. zkouším 'slova'? ne, pět písmen. 'slo...'? zkusím 'sloup'.} \\
\texttt{hádání: sloup} \\
            }
        }
    }
    & & \\ \\

    \theutterance \stepcounter{utterance}  
    & & & \multicolumn{2}{p{0.3\linewidth}}{
        \cellcolor[rgb]{0.9,0.9,0.9}{
            \makecell[{{p{\linewidth}}}]{
                \texttt{\tiny{[GM$|$GM]}}
                \texttt{attempts = 3} \\
\texttt{target\_word = sídlo} \\
\texttt{guess = sloup} \\
\texttt{guess\_feedback = s<green> l<yellow> o<yellow> u<red> p<red>} \\
            }
        }
    }
    & & \\ \\

    \theutterance \stepcounter{utterance}  
    & & & \multicolumn{4}{p{0.6\linewidth}}{
        \cellcolor[rgb]{0.9,0.9,0.9}{
            \makecell[{{p{\linewidth}}}]{
                \texttt{\tiny{[P1$\langle$GM]}}
                \texttt{zpětná vazba k hádání: s<zelená> l<žlutá> o<žlutá> u<červená> p<červená>} \\
\\ 
\texttt{Odpovězte pouze v tomto formátu.} \\
\texttt{vysvětlení: podrobnosti} \\
\texttt{hádání: slovo} \\
            }
        }
    }
     \\ \\

    \theutterance \stepcounter{utterance}  
    & \multicolumn{4}{p{0.6\linewidth}}{
        \cellcolor[rgb]{0.561,0.737,0.561}{
            \makecell[{{p{\linewidth}}}]{
                \texttt{\tiny{[P1$\rangle$GM]}}
                \texttt{vysvětlení: 's' je zelené na první pozici. 'o' je žluté (potvrzeno, není na pozici 3, ani 2, 4 z minula). 'l' je žluté (není na pozici 2, ani 5 z minula). Vyloučeno: t, r, m, k, u, p. Zbývá najít slovo s písmeny s, o, l a dvěma dalšími neznámými. Možné pozice: 'o' může být 3. nebo 5.? (2. a 4. vyloučeno, 3. právě vyloučeno), takže 'o' musí být na 5. pozici. 'l' nemůže být 2., 5. (obsazeno 'o'), ani z minula 5., takže 'l' může být 3. nebo 4. Slovo: s \_ \_ \_ o, kde na 3. nebo 4. je 'l'. Zkouším 'sedlo'? Ale 'd' a 'e' neznám. Zkusím 'selio'? Nesmysl. Zkusím 'sallo'? Nesmysl. Zkusím 'solno'? Neobsahuje 'l' dvakrát. Zkusím 'sídlo'? 's' na začátku, 'o' na konci, 'l' uvnitř. Zkusím 'sadlo'? Ne, 'd' neznám. Zkusím 'sýklo'? Ne. Co 'sello'? Dvě 'l'. Zkusím 'sollo'? Ne. Co 'salvo'? 'v' neznám. Co 'salso'? Ne. Zkusím 'sólo' – to je 4 písmena. Zkusím 'sokol' už bylo. Co 'soplo'? Ne. Zkusím 'sídlo' – 's' na 1., 'í' na 2., 'd' na 3., 'l' na 4., 'o' na 5. 'd' nevyloučeno. Zkusím 'sídlo'.} \\
\texttt{h} \\
            }
        }
    }
    & & \\ \\

    \theutterance \stepcounter{utterance}  
    & & & \multicolumn{2}{p{0.3\linewidth}}{
        \cellcolor[rgb]{0.9,0.9,0.9}{
            \makecell[{{p{\linewidth}}}]{
                \texttt{\tiny{[GM$|$GM]}}
                \texttt{Error: The response should contain the 'hádání:' keyword exactly once.} \\
            }
        }
    }
    & & \\ \\

    \theutterance \stepcounter{utterance}  
    & & & \multicolumn{2}{p{0.3\linewidth}}{
        \cellcolor[rgb]{0.9,0.9,0.9}{
            \makecell[{{p{\linewidth}}}]{
                \texttt{\tiny{[GM$|$GM]}}
                \texttt{game\_result = ABORT} \\
            }
        }
    }
    & & \\ \\

\end{supertabular}
}

\end{document}
